% SPDX-FileCopyrightText: Copyright (c) 2024-2026 Yegor Bugayenko
% SPDX-License-Identifier: MIT

\section{Memory}\label{sec:memory}

The \deff{malloc} object is an abstraction of a block of
  random-access memory.
A copy is created through one of its attributes: \deff{of} allocates
  a block of a given size, \deff{empty} allocates an empty block, and
  \deff{for} allocates a block sized to hold a given object.
The allocated block behaves like \deff{bytes}, but additionally has
  \adeff{write}, \adeff{read}, \deff{copy}, and \deff{put} attributes.
In this example, one kilobyte of memory with random data
  is allocated when a copy of \deff{malloc} is created,
  six bytes are written starting at the 200th byte,
  and then a string is read back and \ff{"Hello!"} is printed:

\begin{ffcode}
malloc.of
  1024
  seq * > [m]
    m.write 200 "Hello!"
    io.stdout
      string
        m.read 200 6
\end{ffcode}

The \deff{malloc} object also supports resizing via
  the \adeff{resized} attribute and copying the data within
  the allocated memory block via the \deff{copy} attribute.
In the next example, \deff{malloc} is allocated and filled with 5 bytes,
  then resized to 10 bytes and its content is duplicated.
The result block contains \ff{"hellohello"}.

\begin{ffcode}
malloc.of
  5
  seq * > [m]
    m.write 0 "hello"
    m.resized 10
    m.copy 0 5 5
\end{ffcode}

The result of dataization of \deff{malloc} is the result
  of dataization of its second argument, which is the scope
  where the memory block is allocated and cleared automatically
  when dataization has happened.
